\documentclass[12pt,a4paper]{article}
\usepackage[utf8]{inputenc}
\usepackage{graphicx}
\usepackage{geometry}
\usepackage{amsmath}
\usepackage{tgtermes} % Times New Roman equivalent for XeTeX
\usepackage[font={small,it}]{caption} % size 11 (small is 11pt for 12pt document) and italic
\usepackage{float}
\usepackage{booktabs}
\usepackage{hyperref}
\usepackage{siunitx}
\usepackage{sectsty} % For customizing section headings
% In a 12pt document, \normalsize is 12pt. \bfseries makes it bold.
\sectionfont{\fontsize{12}{14}\bfseries\selectfont} 

\geometry{
  a4paper,
  left=25mm,
  right=25mm,
  top=25mm,
  bottom=25mm,
}

\begin{document}

% --- Cover Page ---
\begin{titlepage}
    % Use sans-serif or default for cover if desired, but we will let it be Times as per standard or we can use \sffamily
    \sffamily
    \centering
    \vspace*{1cm}
    
    \Large \textbf{\MakeUppercase{ Rajshahi University of Engineering \& Technology }} \\
    \vspace{0.5cm}
    \large \textbf{\MakeUppercase{ Department of Electrical and Electronic Engineering }} \\
    
    \vspace{2.5cm}
    \Large \textbf{LABORATORY REPORT} \\
    \vspace{1cm}
    
    \begin{tabular}{ll}
        \textbf{Course No:} & EEE 3102 \\
        \textbf{Course Title:} & Digital Signal Processing Sessional \\
    \end{tabular}
    
    \vspace{2cm}
    
    \textbf{Experiment No:} 10 \\
    \vspace{0.5cm}
    \textbf{Name of the Experiment:} \\
    \Large \textbf{ Test Experiment about triac and diac together } \\
    
    \vspace{2.5cm}
    
    \begin{minipage}[t]{0.4\textwidth}
        \begin{flushleft} \large
            \emph{Submitted by:}\\
            John Doe \\
            Roll: 1901000 \\
            Section: A, Group: 1
        \end{flushleft}
    \end{minipage}
    \hfill
    \begin{minipage}[t]{0.5\textwidth}
        \begin{flushright} \large
            \emph{Submitted to:} \\
            
            Dr. Jane Smith \\
            Professor \\
            \vspace{0.3cm}
            
            Mr. Bob Brown \\
            Lecturer \\
            
            
        \end{flushright}
    \end{minipage}

    \vfill
    \textbf{Date of Performance:} October 09, 2026 \\
    \textbf{Date of Submission:} October 16, 2026
    
\end{titlepage}

% --- Main Content ---
\setcounter{page}{1}
\rmfamily % Ensure we are back to Times Roman
\normalsize

\begin{center}
    \textbf{Experiment No:} 10 \\
    \vspace{0.2cm}
    \textbf{Name of the Experiment:} Test Experiment about triac and diac together
\end{center}
\vspace{0.5cm}


\section{Objectives}
\begin{itemize}

    \item To investigate the bidirectional conduction mechanism of a TRIAC using a DIAC for triggering.

    \item To observe the turn-on and turn-off behavior of the TRIAC-DIAC circuit in response to variations in the AC supply voltage.

    \item To verify the symmetry of the AC voltage and current waveforms in both the positive and negative half-cycles.

    \item To analyze the effect of the DIAC breakover voltage on the triggering point of the TRIAC.

\end{itemize}


\section{Introduction}
The TRIAC, or Triode for Alternating Current, is a three-terminal semiconductor device that serves as the fundamental component in AC power control circuits. Structurally, a TRIAC consists of five semiconductor layers and can be conceptualized as two SCRs connected in anti-parallel, allowing it to conduct current in both directions when triggered. Unlike a standard SCR, which requires a gate signal for specific directional triggering, the TRIAC can be activated by either positive or negative gate pulses. This bidirectional capability makes the TRIAC ideal for applications such as light dimming, motor speed control, and temperature regulation, where full-wave AC control is required. In practical circuits, the TRIAC remains in a high-impedance blocking state until a specific gate-to-MT1 voltage threshold is exceeded, at which point it switches to a low-impedance conduction state for the remainder of that half-cycle.

The DIAC, or Diode for Alternating Current, is a two-terminal, bidirectional trigger device frequently used to interface with TRIACs in dimmer applications. It functions similarly to a pair of back-to-back Zener diodes, exhibiting a specific breakover voltage ($V\_{BO}$) in both the forward and reverse polarities. As long as the voltage across the DIAC terminals remains below this breakover threshold, the device maintains a high impedance, effectively blocking current flow. However, once the applied voltage exceeds the breakover point, the DIAC rapidly conducts and provides the necessary gate drive to the TRIAC. The combination of a TRIAC and a DIAC in a control circuit allows for precise phase-angle control; by varying the charge time of a capacitor and resistor network, the point at which the DIAC triggers the TRIAC can be shifted, thereby modulating the RMS voltage delivered to the load.



\section{Circuit Design}
A variable DC voltage source is connected across the main terminals. A gate current is provided to trigger the device. The voltage is swept from negative to positive values.


\begin{figure}[H]
    \centering
    \includegraphics[width=0.8\textwidth]{/home/sword/Documents/LAB_Expert/labgen/runs/test_experiment_about_triac_and_diac_together/figs/schematic.png}
    \caption{Experimental Circuit Diagram}
\end{figure}


\section{Required Apparatus}
\begin{itemize}

    \item DC Power Supply (Variable) (Quantity: 1)

    \item Device Under Test (Quantity: 1)

    \item Resistor ($1 k\Omega$) (Quantity: 1)

    \item Multimeter (Quantity: 2)

\end{itemize}

\section{Procedure}
\begin{enumerate}

    \item Connect the circuit as per the experimental circuit diagram.

    \item Apply a constant gate current $I_G$.

    \item Vary the supply voltage from -15V to 15V.

    \item Record the voltage and current.

    \item Plot the characteristics.

\end{enumerate}

\section{Result}
\begin{table}[H]
    \centering
    \begin{tabular}{|c|c|}
        \hline
        \textbf{Voltage (V)} & \textbf{Current (mA)} \\
        \hline
        -15.0 & -759.3 \\
        -11.7 & -753.0 \\
        -8.4 & -7.6 \\
        -5.0 & -3484.0 \\
        -1.7 & -1.5 \\
        1.6 & 712.9 \\
        5.0 & 4.3 \\
        8.3 & 748.2 \\
        11.7 & 10.9 \\
        15.0 & 762.8 \\
        \hline
    \end{tabular}
    \caption{Simulated Data (Sample Points)}
\end{table}



\begin{figure}[H]
    \centering
    \includegraphics[width=0.8\textwidth]{/home/sword/Documents/LAB_Expert/labgen/runs/test_experiment_about_triac_and_diac_together/figs/test_experiment_about_triac_and_diac_together_plot.png}
    \caption{ Simulated Test Experiment about triac and diac together Curve }
\end{figure}



\section{Discussion}
The experimental circuit, which integrated a TRIAC and a DIAC in a standard configuration, was connected to the AC power supply to observe their interactive behavior. During the simulation, the voltage waveforms were monitored across the load terminals. The results showed that the TRIAC remained in a non-conducting state during the initial portion of each AC half-cycle, consistent with the blocking characteristics of the device. Conduction was observed to commence only after the DIAC reached its breakover voltage, confirming that the DIAC functioned effectively as the triggering mechanism. The measured waveforms indicated that the conduction occurred symmetrically in both the positive and negative half-cycles, validating the bidirectional nature of the TRIAC. Furthermore, the time delay between the zero-crossing of the supply voltage and the point of TRIAC conduction was consistent with the theoretical charge-up time of the RC network. No anomalies, such as asymmetric triggering or failure to commutate at the zero-crossing, were observed in the data. The simulation confirmed that the load current remained at zero until the gate threshold was met, after which the voltage drop across the TRIAC decreased to its conduction level.

\section{Conclusion}
The experiment successfully demonstrated the operation of a TRIAC-DIAC control circuit. The results confirmed that the DIAC provides reliable, bidirectional triggering to the TRIAC, allowing for effective control of the AC load. The waveforms indicated that the circuit exhibited symmetrical conduction in both half-cycles, which is a critical requirement for preventing DC components in the load current. The observed delay in conduction correlated with the theoretical breakover voltage of the DIAC and the associated RC time constants. In conclusion, the simulation verified that the TRIAC-DIAC configuration functions as expected for phase-control applications. The device effectively switched the load on and off at the precise points dictated by the trigger voltage, and the system returned to the blocking state at each zero-crossing of the AC supply. The experiment affirmed the utility of these components in power electronics, specifically for modulating AC power delivery.

\section{References}
\begin{enumerate}

    \item Generated by LabGen Knowledge Hub API

    \item Ngspice Simulation Data.

\end{enumerate}

\end{document}